\documentclass[11pt,a4paper]{article}
\usepackage[utf8]{inputenc}
\usepackage[T1]{fontenc}
\usepackage{amsfonts}
\usepackage[french]{babel}
\frenchbsetup{StandardLists=true} % à inclure si on utilise \usepackage[francais]{babel}
\usepackage{enumitem}
\usepackage{amssymb}
\usepackage{amsmath}
\usepackage[left=2cm,right=2cm,top=2cm,bottom=2cm]{geometry}
\usepackage{multirow}
\usepackage[dvipsnames]{xcolor}

\usepackage[T1]{fontenc}
\usepackage{fouriernc}
\usepackage{graphicx}

\usepackage{setspace}
\singlespacing

\usepackage{pstricks}
\usepackage{xkeyval}
\usepackage{pst-labo}
\usepackage{chemist} %chimie


\usepackage{multicol}
\usepackage{caption}
\usepackage{fancyhdr}
\pagestyle{fancy}
\renewcommand\headrulewidth{1pt}
\fancyhead[L]{Lycée Denis Diderot }
\fancyhead[R]{Classe de 1 Bac Pro }
\setcounter{page}{1}\fancyfoot[C]{page \thepage}
\fancyfoot[R]{Année 2024-2025}
\fancyfoot[L]{Alexis RUIZ}
\usepackage{multirow,tabularx,booktabs}
\usepackage[tikz]{bclogo} 

\newcommand{\cadreblanc}[1]{%
\medskip\framebox[\linewidth][c]{%
\rule{0mm}{#1mm}}\par
\medskip}

\usepackage{awesomebox}
\setlength{\aweboxvskip}{0mm}
\usepackage{pifont}
\usepackage{chemmacros}
\chemsetup{modules=all}
\setlength{\columnseprule}{.25mm}

\renewcommand{\arraystretch}{2.5}
\newcommand*\acocher{\quitvmode{\fboxsep0pt \fboxrule0.8pt \fbox{\vrule height1.8ex depth.1ex width0pt \vrule height0pt depth0pt width1.9ex }}}

\newcolumntype{C}[1]{>{\centering\arraybackslash }b{#1}}

\newcommand{\code}{\awesomebox[black]{2pt}{

\includegraphics[scale=0.15]{Chapitre 1 - Energie et puissance électrique/images/frame.png}
}{black}}

\newcommand{\moteur}{\awesomebox[black]{2pt}{

\includegraphics[scale=1]{images/moteur.jpg} 
}{black}}

\newcommand{\defconclusion}{\awesomebox[black]{2pt}{\rotatebox{90}{\Large{\textbf{Conclusion}}}}{black}}

\newcommand{\defbox}{\awesomebox[black]{2pt}{
\includegraphics[scale=0.5]{Chapitre 1 - Energie et puissance électrique/images/appel exam.png}}{black}}

  \newenvironment{itemize*}%  % listes compactes
    {\begin{itemize}%
      \setlength{\itemsep}{0pt}%
      \setlength{\parskip}{0pt}}%
    {\end{itemize}}

\frenchbsetup{StandardLists=true}

%%%%%%%%%%%%%%%%Correction 

\newcommand{\ctext}[1]
{\begin{bclogo}[couleur = white,logo= , arrondi = 0.1,barre=none]
{}
\vspace{-8mm}
~~
\textcolor{red}{#1}
~~
\end{bclogo}}



\begin{document}

\flushleft \textbf{NOM : \hfill Classe : \hspace{3cm} ~~\\[2mm]
Prénom:\\[0,4cm]
}
\hrule\
\\[0,2cm]

\begin{center}
 \LARGE{TP - Préparation d'une solution d'hydroxyde sodium    }
\end{center}

\vspace{3 mm}

\normalsize

\hrule\
\\[0,2cm]
   \begin{center}
   \textbf{Le soin et la rédaction interviendront dans la notation.}
   \end{center} 

   Vous allez préparer 100 mL d'une solution de soude dont la concentration en ions hydroxyde  \ch{OH^-} est \ch{C= 0,2 mol/L} . 
   
\noindent%
\begin{minipage}{.45\textwidth}%

    
 \begin{bclogo}[couleur=white,  arrondi =0.1 , logo = \bcinfo, epBarre =2]{ Matériel }

\begin{itemize*}
    \item 1 balance
    \item 3 béchers
    \item 2 fioles jaugées de 100 mL
    \item 1 entonnoir
    \item 1 pipette avec son dispositif d'aspiration
\end{itemize*}
\end{bclogo}




\end{minipage}%
\hfill
\begin{minipage}{.45\textwidth}%
    
 \begin{bclogo}[couleur=white,  arrondi =0.1 , logo = \bcinfo, epBarre =2]{ Produits }

\begin{itemize*}
    \item Pastilles de soude (formule NaOH)
    \item Eau distillée

\end{itemize*}
\end{bclogo}
%\fbox{\includegraphics[width=\textwidth]{Chapitre 4 - Solutions aqueuses/images/TP2.png}}
\end{minipage}%



\fcolorbox{black}{blue!30}{\textbf{Calculs préparatoires}}
\begin{itemize*}
    \item Calculez la masse m, en g, de soude nécessaire. 
    \ctext{Calcul de la masse molaire M(NaOH) = M(Na) +M(O)+M(H) = 23 + 16 + 1 = 40 g/mol \\[2mm]
    On veut prépare V = 100 mL soit V = 0,1 L de solution. \\[2mm]
    Il nous faut alors n = $C \times V = 0,2 \times 0,1 = 0,02 mol$ \\[2mm]
    On doit donc peser $m= n \times M = 0,02 \times 40 = 0,8 g$
    }

        \item  Les balances du laboratoire sont précises à 0,1 g près, peut-on les utiliser ? justifier votre réponse. 
    \ctext{La balance n'a pas une précision assez importante pour peser 0,8g \\[2mm] 
    En effet avec une précision de 0,1, on aura 0,1/0,8 = 12,5 \% d'erreur possible. \\[2mm]
    Nore masse masse serait comprise entre 0,7 g < m < 0,9 g}
    
    \defbox{\fbox{
\begin{minipage}{8cm}{

Appeler M Ruiz pour qu'il vérifie vos calcul.}
\end{minipage}
}}
%%%%%%%%%%%%%%%%%%%%%%%%%%%%%%%%%
\newpage
%%%%%%%%%%%%%%%%%%%%%%%%%%%%%%%%%%
\noindent%

\begin{minipage}{.90\textwidth}%
\fcolorbox{black}{blue!30}{\textbf{Protocole dissolution}}

\begin{itemize}[label=$\bullet$]
    \item Dans une coupelle, pesez 4,0g de soude (Hydroxyde de sodium) 
\begin{center}
    \includegraphics[scale=0.5]{Chapitre 4 - Solutions aqueuses/images/dilution.png}
\end{center}


    \item Calculer, en mol/L, la concentration molaire \ch{C_0} de cette solution. 

    \ctext{On a pesé 4g, on a donc $n=\dfrac{4}{40} = 0,1 mol$ d'hydroxyde de sodium NaOH. \\[2mm]
    On les a dissout dans V= 100 mL d'eau soit V = 0,1 L.\\[2mm]
    on a alors une concentration $C_0=\dfrac{n}{V}=\dfrac{0,1}{0,1} = 1 mol/L$} 
\end{itemize}
\end{minipage}

\vfill 

  \defbox{\fbox{
\begin{minipage}{12cm}{

Appeler M Ruiz pour qu'il vérifie votre solution et votre calcul}
\end{minipage}
}}

\vfill 

\fcolorbox{black}{blue!30}{\textbf{Protocole dilution}} \\[3mm]

\item Pour réaliser une dilution on utilisera le protocole ci-dessous : 
\begin{center}
    \includegraphics[scale=0.6]{Chapitre 4 - Solutions aqueuses/images/dilution2.png}
\end{center}


\item Réaliser la dilution , on prélèvera 20 mL de solution mère en deux fois.

\item Calculer, en mol/l, la concentration \ch{C} en soude de cette solution diluée. 

\ctext{
Nous avez prélevé 20 mL que nous avons dilué dans 100 mL ; $\dfrac{100}{5}=5$ \\[2mm]
Nous avons dilué 5 fois \\[2mm]
la concentration est donc divisée par 5. $C = \dfrac{C_0}{5} = \dfrac{1}{5} = 0,2 mol/L$ 
}

\vfill 



\end{itemize*}


\end{document}